% Lösungen Einheit 3 von 3 – Einheit 3 (zusammenbau v0.1)
\einheitenkopf{Einheit 3 von 3 · Einheit 3}

\erg{12}{a) Umkehraufgabe: ein Parameter oder eine Grenze ist gesucht \quad b) $\mu = 40$, $\sigma = 3$ (aus der Klammer) \quad c) $\frac{10 - \mu}{2} \approx 1{,}645$ ergibt $\mu \approx 6{,}71$ min (nicht $10$); damit $P(W < 5) \approx 0{,}1963$. \quad d) $c \approx 38\,592$ km (Quantil von A); für Reifen B ist $P(L > c) \approx 0{,}872$ – nein, nur etwa $87{,}2\,\%$ (entscheidend ist die zweite Verteilung). \quad e) Aus I: $\frac{494 - 500}{\sigma} \le -1{,}645$, also $\sigma_{\text{max}} = 3{,}65$. Für diesen Wert ist $P(X \le 488) \approx 0{,}0005 \le 0{,}001$; für kleineres $\sigma$ rückt die Masse näher an $\mu$ und $P(X \le 488)$ wird noch kleiner (Monotonie in $\sigma$ – ein Einzelfall genügt nicht). Also folgt II aus I.}
\erg{13}{a) $P(X < 39) \approx 0{,}0228$ bei $\sigma = 0{,}5$; halbiert $\approx 0{,}0114$. Aus $P(X < 39) = 0{,}0114$ bei $\mu = 40$: $\frac{39 - 40}{\sigma} \approx -2{,}28$, also $\sigma \approx 0{,}44$ mm (nicht $\sigma$ halbieren).}
\erg{14}{a) Fehler: Im Nenner des Exponenten steht $2\sigma^2$, nicht $\sigma$. Richtig: $2\sigma^2 = 18$, also $\sigma = 3$; $\mu = 30$ stimmt. \quad b) $F(\mu)$ ist die Fläche unter der Dichte links von $\mu$. Die Dichte ist symmetrisch um $\mu$, also liegt links genau die Hälfte der Gesamtfläche eins.}
